\documentclass[11pt]{article}
\usepackage[margin=1in]{geometry}
\usepackage{amsmath,amssymb,amsthm,mathtools,booktabs,array,microtype}
\usepackage[hidelinks]{hyperref}
\newtheorem{theorem}{Theorem}[section]
\newtheorem{lemma}[theorem]{Lemma}
\newtheorem{proposition}[theorem]{Proposition}
\newtheorem{corollary}[theorem]{Corollary}
\theoremstyle{remark}\newtheorem{remark}[theorem]{Remark}
\newcommand{\norm}[1]{\left\|#1\right\|_\infty}
\newcommand{\R}{\mathbb R}
\newcommand{\Q}{\mathbb Q}
\newcommand{\eps}{\varepsilon}
\newcommand{\dist}{\operatorname{dist}}
\title{Self-similar counterexamples to the consecutive-degree\newline conjectures of Lorentz}
\author{}
\date{October 28, 2026}
\begin{document}
\maketitle
\begin{abstract}
We construct two continuous functions on $[-1,1]$. The first vanishes at
zero, is not odd, and has identical best uniform polynomial approximants
at every consecutive pair of degrees $(2k+1,2k+2)$. The second is not even
and has identical best approximants at every pair $(2k,2k+1)$. These give
counterexamples to Conjectures $3'$ and $4$ in the formulation of Saff and
Varga. Both functions are uniformly convergent self-similar series with
contraction $-3/5$ and explicit algebraic polynomial forcing. Three and four
finite alternation certificates, respectively, propagate to every degree.
A quantitative stability argument proves the certificates on the entire
interval with the explicit perturbation parameter $10^{-10}$. The finite
algebraic identities and inequalities are verified by independent programs
using exact rational arithmetic. The new constructions are not asserted
here to have been formalized in Lean.
\end{abstract}

\section{The statements and the proof strategy}

For $f\in C[-1,1]$, write $P_n f$ for its unique best uniform approximant
among real polynomials of degree at most $n$, and let
$E_n(f)=\norm{f-P_n f}$. Throughout, $T_k$ is the Chebyshev polynomial
defined by $T_k(\cos\theta)=\cos(k\theta)$.

The two assertions considered here are the normalized oddness conjecture
and the evenness conjecture from \cite{SaffVarga}:
\begin{align*}
 f(0)=0,\quad P_{2k+1}f=P_{2k+2}f\ (k\ge0)
 &\quad\Longrightarrow\quad f\text{ is odd},\\
 P_{2k}f=P_{2k+1}f\ (k\ge0)
 &\quad\Longrightarrow\quad f\text{ is even}.
\end{align*}

\begin{theorem}\label{thm:main}
There exist $f_{\mathrm{o}},f_{\mathrm{e}}\in C[-1,1]$ such that
\begin{align*}
 f_{\mathrm{o}}(0)&=0,&
 P_{2k+1}f_{\mathrm{o}}&=P_{2k+2}f_{\mathrm{o}}\quad(k\ge0),&
 f_{\mathrm{o}}(1)+f_{\mathrm{o}}(-1)&\ne0,\\
 &&P_{2k}f_{\mathrm{e}}&=P_{2k+1}f_{\mathrm{e}}\quad(k\ge0),&
 f_{\mathrm{e}}(1)-f_{\mathrm{e}}(-1)&\ne0.
\end{align*}
Both functions are specified below by algebraic polynomial forcing and the
fixed parameter $\eps=10^{-10}$.
\end{theorem}

The proof has three dependencies. First, Chebyshev composition reduces all
degree conditions to finitely many initial ones. Second, an explicit
baseline has known extremal points and sufficiently strong local gaps to
preserve finite contact inequalities under perturbation. Third, exact
algebraic certificates supply the required initial approximants. We prove
these ingredients in that order and then assemble the two constructions.

\section{Self-similarity and the finite degree reduction}

Set
\[
 a=\frac35,\qquad u=\frac{\sqrt2}{2},\qquad
 \mathcal S h(x)=\sum_{j=0}^{\infty}(-a)^j h(T_{3^j}(x)),
 \qquad G=\mathcal S T_1.
\]
For continuous $h$, the series converges uniformly on $[-1,1]$, since
$\norm{h\circ T_{3^j}}=\norm h$ and $\sum a^j<\infty$. In particular,
$\mathcal S h$ is continuous and
\begin{equation}\label{eq:self}
 \mathcal S h=h-a(\mathcal S h)\circ T_3.
\end{equation}
Commutation of Chebyshev polynomials gives
$\mathcal S T_d=G\circ T_d$.

\begin{lemma}[Chebyshev composition]\label{lem:composition}
For $F\in C[-1,1]$, $m\ge1$, and $n\ge0$,
\[
 P_n(F\circ T_m)=(P_{\lfloor n/m\rfloor}F)\circ T_m.
\]
\end{lemma}
\begin{proof}
Write $x=\cos\theta$. Average any polynomial competitor $p(\cos\theta)$
over the shifts $\theta\mapsto\theta+2\pi j/m$, $0\le j<m$. The target
$F(\cos m\theta)$ is invariant under these shifts, so this averaging
cannot increase the uniform error. Only cosine frequencies divisible by
$m$ survive. The average is therefore $r(\cos m\theta)$ for a polynomial
$r$ of degree at most $\lfloor n/m\rfloor$. Since $T_m$ maps $[-1,1]$
onto itself, its error equals $\norm{F-r}$. This gives equality of the two
best errors, and uniqueness gives the claimed identity.
\end{proof}

\begin{lemma}[Recurrence for best approximants]\label{lem:recurrence}
If $q$ is a polynomial and $f=\mathcal S q$, then
\begin{equation}\label{eq:bestrec}
 P_n f=q-a(P_{\lfloor n/3\rfloor}f)\circ T_3,
 \qquad n\ge\deg q.
\end{equation}
\end{lemma}
\begin{proof}
For $n\ge\deg q$, translation by $q$ preserves the space of degree-at-most
$n$ polynomials. Apply this translation, homogeneity of best approximation,
\eqref{eq:self}, and Lemma~\ref{lem:composition}.
\end{proof}

\begin{proposition}[Finite conditions for the two parity patterns]
\label{prop:finite}
Let $f=\mathcal S q$.
\begin{enumerate}
\item If $\deg q\le7$ and
$P_1f=P_2f$, $P_3f=P_4f$, and $P_5f=P_6f$, then
$P_{2k+1}f=P_{2k+2}f$ for every $k\ge0$.
\item If $\deg q\le8$ and
$P_0f=P_1f$, $P_2f=P_3f$, $P_4f=P_5f$, and $P_6f=P_7f$, then
$P_{2k}f=P_{2k+1}f$ for every $k\ge0$.
\end{enumerate}
\end{proposition}
\begin{proof}
For the first claim, use strong induction on odd $n$. The cases $n=1,3,5$
are assumptions. For odd $n\ge7$, the two floor indices in
\eqref{eq:bestrec} for $n,n+1$ agree unless $n=6r+5$. In that case they
are $2r+1,2r+2$, a smaller odd/even pair covered by induction.

For the second claim, use strong induction on even $n$. The cases
$n=0,2,4,6$ are assumptions. For even $n\ge8$, the floor indices agree
unless $n=6r+2$, in which case they are the smaller even/odd pair
$2r,2r+1$. The same recurrence completes the proof.
\end{proof}

\section{The baseline and its exact extrema}

For $0<c<1$, define
\[
 G_c(x)=\sum_{j\ge0}(-c)^j T_{3^j}(x).
\]
Thus $G=G_a$. Every $G_c$ is odd. We first establish the extremal geometry
for $b=1/3$ and then transfer it to $a=3/5$.

\begin{lemma}\label{lem:baseline}
The function $G$ has unique maximum at $u$, unique minimum at $-u$, and
\begin{equation}\label{eq:B}
 B:=G(u)=\norm G=\frac{u}{1-a}=\frac{5u}{2}.
\end{equation}
Moreover, for every $y\in[-1,1]$,
\begin{equation}\label{eq:globalgap}
 B-|G(y)|\ge\frac12\dist(y,\{-u,u\})^2.
\end{equation}
\end{lemma}
\begin{proof}
Pair adjacent terms in $G_b$:
\[
 G_b=\sum_{j\ge0}9^{-j}g\circ T_{9^j},
 \qquad g(x)=x-\frac13T_3(x)=2x-\frac43x^3.
\]
The identity
\[
 g(u)-g(y)=\frac43(y-u)^2(y+2u)
\]
shows that $g$ has its unique maximum at $u$ and its unique minimum at
$-u$ on the interval. Since $T_9(u)=u$, every term in the paired series
has its maximum at $u$. Consequently $B_b:=\norm{G_b}=3u/2$ and
\begin{equation}\label{eq:bquad}
 B_b-G_b(y)\ge\frac43(y-u)^2(y+2u)\ge\frac12(y-u)^2.
\end{equation}
The last inequality uses $\frac43(\sqrt2-1)>1/2$. Oddness gives the
corresponding lower bound for $B_b+G_b(y)$.

The uniformly convergent series satisfy the identity
\begin{equation}\label{eq:resolvent}
 G_a(x)=G_b(x)+(a-b)\sum_{j=1}^{\infty}
 a^{j-1}(-1)^jG_b(T_{3^j}(x)).
\end{equation}
Indeed, after expansion, the coefficient of $T_{3^k}$ on the right is
\[
 (-1)^k\left[b^k+(a-b)\sum_{j=1}^k a^{j-1}b^{k-j}\right]=(-a)^k.
\]
Rearrangement is justified by absolute uniform convergence.

Every coefficient multiplying a function on the right of
\eqref{eq:resolvent} is positive. Each such function has norm $B_b$ and
value $B_b$ at $u$, because $T_3(u)=-u$ and $G_b$ is odd. Thus
\[
 G_a(u)=B_b\left(1+\frac{a-b}{1-a}\right)=\frac{u}{1-a}.
\]
All deficits from this maximum are nonnegative. Keeping the first deficit
and using \eqref{eq:bquad} proves
$B-G_a(y)\ge(y-u)^2/2$. Oddness proves
$B+G_a(y)\ge(y+u)^2/2$. Their minimum is
\eqref{eq:globalgap}, and also proves uniqueness of the two extrema.
\end{proof}

For later use, $P_0G=0$ by oddness. Lemma~\ref{lem:recurrence} applied to
$G=\mathcal S T_1$ therefore gives
\begin{equation}\label{eq:Gsmall}
 P_1G=P_2G=T_1.
\end{equation}
It follows that $G\circ T_d$ has initial best approximant zero for
$0\le n<d$, and best approximant $T_d$ for $d\le n<3d$. The residuals
in these two ranges are $G\circ T_d$ and $-aG\circ T_{3d}$.

\section{From finite contact inequalities to the whole interval}

The point of the larger contraction is that the baseline gap and an
arbitrary polynomial-forcing perturbation can be compared at the same
scale. No stationarity conditions on the perturbation are needed.

\begin{lemma}[A gap at the exit scale]\label{lem:exit}
Let $0<|s|\le\delta:=1/16$, and choose $J\ge0$ minimal with
$3^J|s|\ge\delta$. Then
\begin{equation}\label{eq:exitgap}
 B-G(\cos(\pi/4+s))\ge\frac{a^J}{4608}.
\end{equation}
\end{lemma}
\begin{proof}
Minimality gives $\delta\le3^J|s|\le3\delta$. The exact phase relation,
using $3^J\equiv1$ or $3\pmod8$ according to the parity of $J$, is
\[
 (-1)^JG_b(\cos(3^J(\pi/4+s)))
 =G_b(\cos(\pi/4+(-1)^J3^Js)).
\]
On $[\pi/4-3\delta,\pi/4+3\delta]$, sine is at least $1/2$.
Thus the argument of the final $G_b$ is at distance at least
$\delta/2$ from $u$, and its deficit is at least
$\delta^2/8=1/2048$ by \eqref{eq:bquad}.
For $J\ge1$ the coefficient of this term in \eqref{eq:resolvent} is
$(a-b)a^{J-1}=(4/9)a^J$. For $J=0$ the coefficient is $1\ge4/9$.
Discarding all other nonnegative deficits proves \eqref{eq:exitgap}.
\end{proof}

For a polynomial $w=\sum w_kT_k$, the bounds
\[
 \norm w\le\sum|w_k|,\qquad
 \sup_\theta\left|\frac{d}{d\theta}w(\cos\theta)\right|
 \le\sum k|w_k|
\]
follow directly from the cosine representation. The second bound is in
the angular variable, not the spatial variable.

\begin{lemma}[Perturbation oscillation]\label{lem:oscillation}
Let $H=\mathcal S h$ and $D=H-v$, where $h,v$ are polynomials. Suppose
\[
 \norm h\le M,\quad
 \sup_\theta|(h(\cos\theta))'|\le L,\quad
 \norm v\le V,\quad
 \sup_\theta|(v(\cos\theta))'|\le L_v.
\]
Let $d\ge1$ and let $J$ be a nonnegative integer with
$\delta\le3^Jd|t|\le3\delta$. Then
\begin{equation}\label{eq:oscillation}
 |D(\cos(\theta+t))-D(\cos\theta)|\le C_d a^J,
 \qquad
 C_d=5M+\frac{15L}{64d}+\frac{3L_v}{16d}.
\end{equation}
Also $\norm D\le K:=\frac52M+V$.
\end{lemma}
\begin{proof}
The $j$th forcing difference is at most
$\min\{L3^j|t|,2M\}$. Splitting the series before $J$ gives
\begin{align*}
 |H(\cos(\theta+t))-H(\cos\theta)|
 &\le L|t|\sum_{j<J}(3a)^j+2M\sum_{j\ge J}a^j\\
 &\le\left(\frac{15L}{64d}+5M\right)a^J.
\end{align*}
Here $3a-1=4/5$ and $3^J|t|\le3/(16d)$.
Moreover $3^{-J}\le a^J$, so
$L_v|t|\le3L_va^J/(16d)$. This proves \eqref{eq:oscillation}.
The norm bound follows from $\norm H\le M/(1-a)=5M/2$.
\end{proof}

\begin{proposition}[Quantitative finite-contact stability]\label{prop:stability}
Let
\[
 R(x)=sA G(T_d(x)),\qquad A>0,\quad s\in\{-1,1\},\quad d\ge1.
\]
Its extremal points are the $2d$ zeros of $T_{2d}$, namely
$x_i=\cos((2i+1)\pi/(4d))$, $0\le i<2d$. Put
$\sigma_i=R(x_i)/(AB)$. For $D=\mathcal S h-v$, suppose
\begin{equation}\label{eq:finiteineq}
 \sigma_iD(x_i)\le\lambda\qquad(0\le i<2d).
\end{equation}
With the constants of Lemma~\ref{lem:oscillation}, set
\[
 \eta_0=\frac1{1296},\qquad
 \gamma=\frac1{8\cdot1296^2}.
\]
If $\eps>0$ satisfies
\begin{equation}\label{eq:epsconditions}
 \eps C_d\le\frac{A}{4608},\qquad
 \eps(K+|\lambda|)\le A\gamma,
\end{equation}
then
\begin{equation}\label{eq:globalcertificate}
 |R(x)+\eps D(x)|\le AB+\eps\lambda\qquad(x\in[-1,1]).
\end{equation}
Every equality in \eqref{eq:finiteineq} is an equality in
\eqref{eq:globalcertificate} with sign $\sigma_i$.
\end{proposition}
\begin{proof}
Use $x=\cos\theta$. Around every extremal angle $\theta_i$ take the
neighborhood $|\theta-\theta_i|\le\eta_0/d$. By periodicity, reflection,
and oddness, the signed baseline value reduces to
$AG(\cos(\pi/4+s'))$, with $|s'|=d|\theta-\theta_i|$.
For $s'\ne0$ choose its exit scale $J$. Lemmas~\ref{lem:exit} and
\ref{lem:oscillation} give
\[
 \sigma_i(R(x)+\eps D(x))
 \le AB-\frac{A a^J}{4608}+\eps\lambda+\eps C_d a^J
 \le AB+\eps\lambda.
\]
At the contact itself this follows from \eqref{eq:finiteineq}.

We must also control the opposite sign. The neighborhood radius ensures
$J\ge4$. Apply the forcing estimate in Lemma~\ref{lem:oscillation} to
$G=\mathcal S T_1$, taking $d=1$, $M=L=1$, and omitting $v$. Then
\[
 |G(\cos(\pi/4+s'))-B|
 \le\frac{335}{64}a^J
 \le\frac{335}{64}\left(\frac35\right)^4
 <\frac{679}{1000}<\frac B2.
\]
Thus $\sigma_iR(x)>AB/2$. The second inequality in
\eqref{eq:epsconditions} gives $\eps\norm D\le A\gamma<AB/2$.
The perturbed residual retains this sign, proving the absolute bound
throughout the neighborhoods. At $s'=0$, sign retention follows directly
from $AB-\eps K>0$.

Outside these neighborhoods, the angular distance of $d\theta$ from
every preimage of $\pm u$ is at least $\eta_0$. Monotonicity of cosine
between its extrema, and the sine bound $1/2$ in a short interval adjacent
to each preimage, give
$\dist(\cos(d\theta),\{-u,u\})\ge\eta_0/2$.
By \eqref{eq:globalgap}, $|R(x)|\le AB-A\gamma$. Therefore
\[
 |R(x)+\eps D(x)|\le AB-A\gamma+\eps K\le AB+\eps\lambda.
\]
Finally, $AB+\eps\lambda>0$ because
$\eps|\lambda|\le A\gamma<AB$. At an active contact the signed value
is exactly $AB+\eps\lambda$, as claimed.
\end{proof}

In both constructions below the exact certificates give
\begin{equation}\label{eq:universalbounds}
 M<32,\quad L<128,\quad V<32,\quad L_v<128,\quad|\lambda|<32.
\end{equation}
All residuals have $A\ge3/5$ and $d\ge3$. Thus $C_d<178$ and
$K+|\lambda|<144$. The single explicit choice
\begin{equation}\label{eq:epsilon}
 \eps=10^{-10}
\end{equation}
satisfies \eqref{eq:epsconditions}, since
\[
 \frac{178}{10^{10}}<\frac{3}{5\cdot4608},\qquad
 \frac{144}{10^{10}}<\frac{3}{5\cdot8\cdot1296^2}.
\]
Consequently, any of the finite certificates below with the stated bounds
is automatically a certificate on the whole interval.

\section{The normalized oddness counterexample}

Define
\begin{align}
 h_{\mathrm{o}}={}&(\eta-\sqrt3-1)T_0-\sqrt3\,T_2+T_4-3T_5
                  +\eta T_6+\beta T_7,\label{eq:oddh}\\
 \eta={}&\frac{16\cos(\pi/18)-26\cos(\pi/6)-26\cos(5\pi/18)}{15}.
 \label{eq:eta}
\end{align}
To specify $\beta$ exactly, put $c_j=\cos(j\pi/36)$ and set
\begin{align}
 7132971\beta={}&-21647763+7264892c_1-16364700c_2+2333622c_3
 -20710080c_4\notag\\
 &+7960924c_5-8377290c_6+6354480c_7-6380100c_8\notag\\
 &+8032814c_9-667710c_{10}+24678c_{11}.
 \label{eq:beta}
\end{align}
Numerically $\eta\approx-1.564814286803340$ and
$\beta\approx-5.914876184323304$. These decimals are descriptive only;
\eqref{eq:eta}--\eqref{eq:beta} define the coefficients.

Set
\begin{equation}\label{eq:fo}
 f_{\mathrm{o}}=G\circ T_3+\eps\mathcal S h_{\mathrm{o}}
 =\mathcal S(T_3+\eps h_{\mathrm{o}}).
\end{equation}
We next give the three finite certificates needed for this function.

\subsection{Exact definition of the correction polynomials}

Let
\[
 x_j=-\cos\frac{(2j+1)\pi}{36},\quad0\le j<18,\qquad
 \sigma_j=(1,-1,-1,1)_{j\bmod4}.
\]
The coarse grid is $y_i=x_{3i+1}$, $0\le i<6$, with signs
$\tau=(-1,1,1,-1,-1,1)$. These are respectively the extrema of
$-aG\circ T_9$ and $G\circ T_3$.

Define the explicit grid functions
\[
 W_k(x_j)=T_k(x_j)-aT_{3k}(x_j)
        +\frac{a^2}{1+(-1)^ka}T_{9k}(x_j).
\]
They equal $(\mathcal S T_k)(x_j)$: after two iterates the finite orbit
is on the two-cycle $\{\pm u\}$, and the remaining series is geometric.
Equivalently, the exact identities
$W_k(x_j)+aW_k(T_3x_j)=T_k(x_j)$ determine these values uniquely.
Write $H_j=\sum_{k=0}^7(h_{\mathrm{o}})_k W_k(x_j)$.

The correction polynomials $v_1,v_3,v_5$ and scalars
$\lambda_2,\lambda_4,\lambda_6$ are specified by the contact equations
\begin{align*}
 v_1(y_i)+\tau_i\lambda_2&=H_{3i+1},
       &&i\in I_2:=\{0,1,4,5\},\\
 v_3(x_j)+\sigma_j\lambda_4&=H_j,
       &&j\in I_4:=\{0,2,4,13,15,17\},\\
 v_5(x_j)+\sigma_j\lambda_6&=H_j,
       &&j\in I_6:=\{0,2,4,6,11,13,16,17\},
\end{align*}
with degrees at most $1,3,5$, respectively. The systems are consistent
for the exact values \eqref{eq:eta}--\eqref{eq:beta}. Each has full column
rank, so the correction and scalar are unique. Their exact rational
coordinates in the field described below are included in
\texttt{odd\_certificate.json}.

For orientation, their Chebyshev coefficients are:
\begin{center}\small
\begin{tabular}{c r r r}
\toprule
coefficient & $v_1$ & $v_3$ & $v_5$\\\midrule
$T_0$&$-3.310540684$&$-3.310540684$&$-3.380936697$\\
$T_1$&$-0.781038720$&$ 0.805215314$&$-1.373864474$\\
$T_2$&$0$&$-1.138509891$&$-0.999856812$\\
$T_3$&$0$&$-2.959735280$&$ 1.640309342$\\
$T_4$&$0$&$0$&$-0.070396013$\\
$T_5$&$0$&$0$&$-4.379580520$\\\bottomrule
\end{tabular}
\end{center}
The corresponding scalars are approximately
$4.600544015619$, $7.815043886301$, and $6.118310129767$.

\begin{lemma}[Odd-pattern finite certificate]\label{lem:oddcertificate}
The above corrections satisfy the bounds \eqref{eq:universalbounds}.
At every coarse or fine node, as appropriate,
\[
 \tau_i(H_{3i+1}-v_1(y_i))\le\lambda_2,
 \qquad
 \sigma_j(H_j-v_{n-1}(x_j))\le\lambda_n\quad(n=4,6).
\]
Equality holds exactly on $I_2,I_4,I_6$, respectively. Each active set has
alternating signs. Every inactive slack is greater than $3/50$.
\end{lemma}
\begin{proof}[Exact arithmetic verification]
Put $z=2\cos(\pi/36)$. All entries lie in
\[
 \Q[z],\qquad
 M(z)=z^{12}-12z^{10}+54z^8-112z^6+105z^4-36z^2+1=0.
\]
The relevant root is isolated in $(1.98,2)$. To see uniqueness there, write
$M=S_6^2-3$, where $S_6=z^6-6z^4+9z^2-2$; on that interval
$S_6>0$ and $S_6'=6z(z^2-1)(z^2-3)>0$.

The program \texttt{verify\_odd.py} constructs every coefficient by exact
field arithmetic, checks the contact equations, and checks the displayed
closed expressions for $\eta,\beta$. It encloses the root by 250 rational
bisections and evaluates every nonzero slack by rational interval
arithmetic. The lower bounds exceed $3/50$; the least positive slacks for
the three certificates exceed $1.5265$, $1.6678$, and $0.0619$,
respectively. Exact interval bounds also establish
\eqref{eq:universalbounds}, $\eta<0$, and positivity of the three scalars.

The independently implemented \texttt{audit\_odd.py} starts from the
serialized rational coefficients. It checks the root isolation, ordered
nodes, cubic node map, finite functional equation for the actual $H$,
all slack identities, their signs, and the alternating active sets. It
also checks \eqref{eq:epsconditions} directly for $\eps=10^{-10}$.
These programs use rational arithmetic for every acceptance condition;
floating-point conversion is used only in printed summaries.
\end{proof}

\subsection{Completion of the first construction}

For degrees $1,2$, the baseline approximant is zero, so the residual for
the candidate $\eps v_1$ is
$G\circ T_3+\eps(\mathcal S h_{\mathrm{o}}-v_1)$.
For degrees $3,4$ and $5,6$, the baseline approximant is $T_3$, and the
residual is $-aG\circ T_9+\eps(\mathcal S h_{\mathrm{o}}-v)$.
Apply Proposition~\ref{prop:stability} with
$(A,s,d)=(1,1,3)$ or $(a,-1,9)$.
The finite certificate provides respectively four, six, and eight
alternating contacts, while the candidates have degrees at most $1,3,5$.
Chebyshev alternation therefore gives
\[
\begin{array}{c|c|c}
 n&P_n f_{\mathrm{o}}&E_n(f_{\mathrm{o}})\\\hline
 1,2&\eps v_1&B+\eps\lambda_2\\
 3,4&T_3+\eps v_3&aB+\eps\lambda_4\\
 5,6&T_3+\eps v_5&aB+\eps\lambda_6.
\end{array}
\]
Proposition~\ref{prop:finite} now proves every odd/even plateau.

The identity $h_{\mathrm{o}}(0)=0$ follows directly from
$T_{2j}(0)=(-1)^j$. Thus \eqref{eq:self} at zero gives
$f_{\mathrm{o}}(0)=0$. At the endpoints,
\[
 h_{\mathrm{o}}(1)+h_{\mathrm{o}}(-1)=4(\eta-\sqrt3).
\]
Since $T_3(\pm1)=\pm1$ and $G\circ T_3$ is odd,
\begin{equation}\label{eq:oddwitness}
 f_{\mathrm{o}}(1)+f_{\mathrm{o}}(-1)
 =\frac{4\eps(\eta-\sqrt3)}{1+a}
 =\frac52\eps(\eta-\sqrt3)\ne0.
\end{equation}
This proves the first assertion of Theorem~\ref{thm:main}.

\section{The evenness counterexample}

This construction uses the baseline $G\circ T_6$ and a degree-eight
forcing perturbation. We specify its coefficients by a small, explicit
linear system, avoiding long rational coordinate lists in the text.

\subsection{Exact algebraic definition}

The unknown forcing is
\[
 h_{\mathrm{e}}=c_1T_1+c_2T_2+c_3T_3+c_4T_4+c_5T_5+c_7T_7+c_8T_8.
\]
Thus $c_0=c_6=0$. Introduce correction polynomials
$w_0,w_2,w_4,w_6$ of degrees at most $0,2,4,6$, and scalars
$\mu_1,\mu_3,\mu_5,\mu_7$.

For $m=12,36$, define the increasing nodes
\[
 x_j^{(m)}=-\cos\frac{(2j+1)\pi}{2m},\qquad0\le j<m.
\]
Use $m_n=12$ for $n=1,3,5$ and $m_7=36$. Let
\[
 \sigma_{n,j}=\begin{cases}
 (1,-1,-1,1)_{j\bmod4},&n=1,3,5,\\
 (-1,1,1,-1)_{j\bmod4},&n=7.
 \end{cases}
\]
For a grid angle $\theta$ set $J=1$ when $m=12$ and $J=2$ when $m=36$,
and define
\begin{equation}\label{eq:evenW}
 W_k^{(m)}(\cos\theta)=
 \sum_{i=0}^{J-1}(-a)^i\cos(k3^i\theta)
 +\frac{\sum_{i=J}^{J+3}(-a)^i\cos(k3^i\theta)}{1-a^4}.
\end{equation}
After $J$ iterates these node orbits have period dividing four. Therefore
\eqref{eq:evenW} is the exact value of $\mathcal S T_k$ at the node.

Write $\mathcal H_j^{(m)}=\sum_k c_k W_k^{(m)}(x_j^{(m)})$.
The defining equations are
\begin{equation}\label{eq:evenlinear}
 w_{n-1}(x_j^{(m_n)})+\sigma_{n,j}\mu_n
 =\mathcal H_j^{(m_n)},\qquad j\in J_n,
\end{equation}
where
\begin{align*}
 J_1&=\{3,6,8\},\\
 J_3&=\{1,3,4,6,8,10\},\\
 J_5&=\{0,2,4,6,7,9,11\},\\
 J_7&=\{3,6,8,13,19,22,27,30,32\}.
\end{align*}
Add the two normalizations
\begin{equation}\label{eq:evennormalize}
 [T_2]w_2=1,\qquad [T_4]w_4=-1.
\end{equation}
Here $[T_k]w$ denotes the coefficient in the Chebyshev basis.
There are $27$ equations for the $7+1+3+5+7+4=27$ unknowns.
Exact elimination shows that this system is nonsingular. Equations
\eqref{eq:evenW}--\eqref{eq:evennormalize} therefore define the forcing
and corrections unambiguously by real algebraic numbers.

Their approximate forcing coefficients are
\[
\begin{array}{c|r}
 k&c_k\\\hline
 1&-0.017561786853260\cr
 2& 0.286725254171375\cr
 3& 0.018613141564351\cr
 4&-0.862027229474975\cr
 5&-0.013885677377745\cr
 7& 0.026299792766218\cr
 8&-0.253316709984462.
\end{array}
\]
The exact rational coordinate lists are in \texttt{even\_certificate.json}.
Define
\begin{equation}\label{eq:fe}
 f_{\mathrm{e}}=G\circ T_6+\eps\mathcal S h_{\mathrm{e}}
 =\mathcal S(T_6+\eps h_{\mathrm{e}}).
\end{equation}

\begin{lemma}[Even-pattern finite certificate]\label{lem:evencertificate}
For each $n=1,3,5,7$ and every point of the corresponding grid,
\[
 \sigma_{n,j}\bigl(\mathcal H_j^{(m_n)}-w_{n-1}(x_j^{(m_n)})\bigr)
 \le\mu_n.
\]
The equality sets are exactly $J_n$. Alternating subsets of cardinality
$n+2$ are $J_1,J_5,J_7$, and
$\{1,3,6,8,10\}\subset J_3$. All nonzero slacks exceed $1/50$.
The coefficient bounds \eqref{eq:universalbounds} hold, and
\begin{equation}\label{eq:oddcoeffsum}
 c_1+c_3+c_5+c_7>\frac1{100}.
\end{equation}
\end{lemma}
\begin{proof}[Exact arithmetic verification]
All coefficients belong to the real algebraic field generated by
$z=2\cos(\pi/72)$, whose defining polynomial is
\begin{align*}
 M_{24}(z)={}&z^{24}-24z^{22}+252z^{20}-1520z^{18}+5814z^{16}
 -14688z^{14}\\
 &+24752z^{12}-27456z^{10}+19305z^8-8008z^6
 +1716z^4-144z^2+1.
\end{align*}
It satisfies $M_{24}(z)=M(z^2-2)$, where $M$ is the polynomial used in
the odd certificate. The relevant root is isolated in $(1.998,2)$;
the same monotonicity argument, applied to $z^2-2$, proves uniqueness.

The generator \texttt{generate\_even.py} solves the $27$-by-$27$ system
by exact arithmetic, with a nonzero pivot at every step. It checks every
functional identity, equality, and slack and encloses the root by $400$
rational bisections. The minimum positive slacks for $n=1,3,5,7$ exceed
$0.1336$, $0.0936$, $0.0965$, and $0.0230$, respectively. Exact coefficient
bounds give $M<2$, $L<34$, $V<2$, $L_v<72$, and $|\mu_n|<1$,
which imply \eqref{eq:universalbounds}; it also verifies
\eqref{eq:oddcoeffsum}.

The independent standard-library program \texttt{audit\_even.py} uses a
separate implementation of rational polynomial arithmetic. From the
serialized coefficients it reconstructs all $36$ fine nodes, verifies
their ordering and cubic map, detects the finite orbit cycles, and
evaluates the exact orbit sums. It then verifies every slack, equality,
alternating set, coefficient bound, and \eqref{eq:epsconditions}. No
floating-point tolerance is used to accept any identity or inequality.
\end{proof}

For reference, the scalar error corrections are approximately
\begin{align*}
 \mu_1&=0.577201789619347,&\quad \mu_3&=0.223648399026073,\\
 \mu_5&=-0.129582723095644,&\quad \mu_7&=0.276968241048616.
\end{align*}
The negative value of $\mu_5$ causes no difficulty: the stability
proposition uses $|\mu_5|$ in its smallness condition, and the total error
remains positive.

\subsection{Completion of the second construction}

By \eqref{eq:Gsmall} and Chebyshev composition, the baseline
$G\circ T_6$ has best approximant zero for degrees $0$ through $5$ and
best approximant $T_6$ for degrees $6,7$. Its residuals are consequently
$G\circ T_6$ and $-aG\circ T_{18}$.

Apply Proposition~\ref{prop:stability} with
$(A,s,d)=(1,1,6)$ for the first three certificates and
$(a,-1,18)$ for the fourth. Alternation gives
\[
\begin{array}{c|c|c}
 n&P_n f_{\mathrm{e}}&E_n(f_{\mathrm{e}})\\\hline
 0,1&\eps w_0&B+\eps\mu_1\\
 2,3&\eps w_2&B+\eps\mu_3\\
 4,5&\eps w_4&B+\eps\mu_5\\
 6,7&T_6+\eps w_6&aB+\eps\mu_7.
\end{array}
\]
Since $\deg(T_6+\eps h_{\mathrm{e}})\le8$,
Proposition~\ref{prop:finite} extends these equalities to every even/odd
pair.

Finally, $G\circ T_6$ is even, while the endpoint difference of the
perturbation is
\begin{equation}\label{eq:evenwitness}
 f_{\mathrm{e}}(1)-f_{\mathrm{e}}(-1)
 =\frac{2\eps}{1+a}(c_1+c_3+c_5+c_7)
 >\frac{\eps}{80}>0.
\end{equation}
This proves the second assertion of Theorem~\ref{thm:main}.

\section{Verification, scope, and reproducibility}

The proof consists of the analytic estimates above and finite exact
algebraic certificates. The latter are not numerical Remez calculations:
all accepted identities are polynomial identities modulo a defining
polynomial, and all accepted strict inequalities are certified by rational
intervals at a uniquely isolated real root. The independently implemented
audits also check that the finite orbit values are the values of the
actual infinite series and that the explicit perturbation size satisfies
the whole-interval estimates.

From the accompanying directory, run
\begin{verbatim}
python3 audit_odd.py
python3 audit_even.py
\end{verbatim}
Both audits use only the Python standard library and the supplied JSON
certificates. To regenerate the odd certificate, run
\texttt{python3 verify\_odd.py}. The optional even generator uses
\texttt{python-flint}, installed by
\texttt{python3 -m pip install python-flint}, and is run with
\texttt{python3 generate\_even.py}. Its output is independently audited
without that dependency.

These constructions are distinct from the earlier counterexample to the
common-origin-zero conjecture. Their best approximants are required to
have the indicated degree plateaus; they are not required to vanish at
zero. The first function itself is normalized at zero. The preceding
Lean formalization of the common-origin-zero example is not a Lean
formalization of the two new results here. The new proof status is a full
analytic argument with two independently checked exact arithmetic
certificates.

\begin{thebibliography}{1}
\bibitem{SaffVarga}
E. B. Saff and R. S. Varga,
\emph{Remarks on some conjectures of G. G. Lorentz},
Journal of Approximation Theory \textbf{30} (1980), 29--36.
\url{https://www.math.kent.edu/~varga/pub/paper_126.pdf}.
\end{thebibliography}
\end{document}
